% Imported from https://github.com/Luca-Yucheng-Jin/QFT-soln
% Source commit: 0bb3e7066154a8fdb256399a8ecceae720c5e192
\section{David Tong The Standard Model, Sheet 3: Anomalies}
\markboth{\thesection\quad TONG THE STANDARD MODEL, SHEET 3}{}

\noindent\textit{These prompts paraphrase David Tong's \emph{The Standard Model},
Example Sheet 3 (February 2026). Problem numbers, starred status, equations,
and guidance follow the
\href{https://davidtong.org/pdfs/teaching/standard-model/rw3.pdf}{official sheet}.}

\subsection{Sheet 3, Problem 1: Chiral $U(1)$ Charge Assignments}

Prove that four left-handed Weyl fermions with integer charges cannot
form a consistent chiral $U(1)$ gauge theory that can also couple to
gravity. Exhibit a consistent example with five such fermions.
\begin{solution}
    There are two constraints in this case,
    \begin{equation}
        \sum_{\rm left}q_i^3 = \sum_{\rm right}q_i^3, \quad \sum q_i = 0.
    \end{equation}
    In the case of four Weyl spinors, we must either have one or two negative charges, namely $\{x,y,z,-w\}$ or $\{x,y,-z,-w\}$ where $x,y,z,w \in \mathbb{Z}^+$. In the first case, we have two equations
    \begin{equation}
        x^3 + y^3+z^3 = w^3, \quad x+y+z=w,
    \end{equation}
    which becomes
    \begin{equation}
        x^2y + xy^2 + y^2z + yz^2 + z^2x + zx^2 = 0,
    \end{equation}
    which only has the trivial solution $x = y =z = w=0$. The same argument works for the second case. Therefore, there cannot be any consistent chiral gauge theory with four left-handed spinors.

    For five left-handed Weyl fermions, $\{1,5,-7,-8,9\}$ does the job.
\end{solution}

\subsection{Sheet 3, Problem 2*: Gauge and 't Hooft Anomalies in a Chiral Theory}

Couple $SU(N)$ to one left-handed Weyl fermion $\chi$ in the
two-index symmetric representation and to $p$ left-handed Weyl
fermions $\psi_i$ in the anti-fundamental. Determine $p$ for gauge
anomaly cancellation.

Classically, $SU(p)$ rotates the $\psi_i$, while $U(1)_\chi$
rotates only $\chi$ and $U(1)_\psi$ rotates only the $\psi_i$.
Show that each Abelian symmetry is anomalous and find a surviving
linear combination. Calculate its $SU(p)^3$, $SU(p)^2U(1)$,
$U(1)^3$, and mixed $U(1)$--gravity 't Hooft anomalies.

The proposed confined spectrum consists of gauge-singlet fermions
\begin{equation}
    \lambda_{ij}=\psi_{[i}(\chi\psi_{j]}),
\end{equation}
in the antisymmetric representation of the global $SU(p)$.
Find their surviving $U(1)$ charge and match all the anomalies
computed above.

\textit{Useful source data:} With $I(\bar R)=I(R)$ and
$A(\bar R)=-A(R)$, the $SU(N)$ representation data are
\begin{center}
\begin{tabular}{c|cccc}
    $R$ & fundamental & adjoint & symmetric & antisymmetric\\ \hline
    $\dim R$ & $N$ & $N^2-1$ & $N(N+1)/2$ & $N(N-1)/2$\\
    $I(R)$ & $1$ & $2N$ & $N+2$ & $N-2$\\
    $A(R)$ & $1$ & $0$ & $N+4$ & $N-4$
\end{tabular}
\end{center}
\begin{solution}
    For the theory to be consistent, we must have
    \begin{equation}
        \sum_i A(R_i) = A(\mathrm{sym}) + pA(\bar{\mathbf{N}}) = N+4 -p = 0 \implies p = N+4.
    \end{equation}

    Then, for the $U(1)_\chi SU(N)^2$ mixed anomaly, we have
    \begin{equation}
        \mathcal{A}_{U(1)_\chi SU(N)^2} = \sum q_i I(R_i) = I(\mathrm{sym}) = N+2,
    \end{equation}
    while for the $U(1)_\psi SU(N)^2$
    \begin{equation}
        \mathcal{A}_{U(1)_\psi SU(N)^2} = \sum q_i I(R_i) = pI(\bar{\mathrm{fund}}) = N+4,
    \end{equation}
    which are both anomalous. For a general $U(1)$ transformation, let us take the charge of $\chi$ to be $Q$ and the charges of $\psi_i$ to be $q$. For the symmetry to be anomaly-free, we must have
    \begin{equation}
        \mathcal{A}_{U(1) SU(N)^2} = pq + Q(N+4) = 0.
    \end{equation}
    A convenient choice of charge is $ q= -(N+2)$ and $Q = N+4$.

    Now, we turn to the computation of the 't Hooft anomalies:
    \begin{equation}
        \mathcal{A}_{SU(p)^3} = \sum_{\text{gauge index}} A(\mathbf{p}) = N,
    \end{equation}
    \begin{equation}
        \mathcal{A}_{SU(p)^2U(1)} = \sum_{\text{gauge index}}q I(\mathbf{p}) = -N(N+2),
    \end{equation}
    \begin{equation}
        \mathcal{A}_{U(1)^3} = \sum_{\text{gauge index}} \sum_i q^3_i = Npq^3 + \dim(\mathrm{sym}) Q^3 = -\frac{N^3(N+3)(N+4)}{2},
    \end{equation}
    \begin{equation}
        \mathcal{A}_{U(1)} = \sum_{\text{gauge index}} \sum_i q_i = Npq + \dim(\mathrm{sym}) Q = -\frac{N(N+4)(N+3)}{2}.
    \end{equation}

    Under $U(1)$ the $\lambda$ fermion transform as
    \begin{equation}
        U(1):\quad \lambda^{ij} \to e^{i(Q+2q)\alpha}\lambda^{ij},
    \end{equation}
    which has a charge $q_\lambda = -N$.

    Since the $\lambda$ fermion is a gauge singlet, the 't Hooft anomalies of this theory are given by
    \begin{equation}
        \mathcal{A}_{SU(p)^3}' =  A(\mathrm{antisym}) = p-4 = N,
    \end{equation}
    \begin{equation}
        \mathcal{A}_{SU(p)^2U(1)}' = q_\lambda I(\mathrm{antisym}) = -N(p-2) = -N(N+2),
    \end{equation}
    \begin{equation}
        \mathcal{A}_{U(1)^3}' =  \sum_{\text{fermion index}} q_\lambda^3 = -\frac{1}{2}p(p-1)N^3 = -\frac{N^3(N+3)(N+4)}{2},
    \end{equation}
    \begin{equation}
        \mathcal{A}_{U(1)}' = \sum_{\text{fermion index}} q_\lambda = -\frac{1}{2}p(p-1)N =  -\frac{N(N+3)(N+4)}{2}.
    \end{equation}
    Therefore, we see that the anomalies indeed match.
\end{solution}

\subsection{Sheet 3, Problem 3*: Real Hypercharges and Anomaly Cancellation}

Use one Standard Model generation, omitting a right-handed neutrino,
with arbitrary real hypercharges:
\begin{equation}
    Q_L:(\mathbf2,\mathbf3)_q,\quad
    L_L:(\mathbf2,\mathbf1)_l,\quad
    u_R:(\mathbf1,\mathbf3)_u,\quad
    d_R:(\mathbf1,\mathbf3)_d,\quad
    e_R:(\mathbf1,\mathbf1)_x.
\end{equation}
The entries specify $SU(2)\times SU(3)$ representations and the
subscript is the $U(1)$ charge.
\begin{enumerate}[label=\roman*)]
    \item Write every gauge anomaly-cancellation condition, including
    the mixed $U(1)$--gravity condition.
    \item Show that the real solutions consist of two branches: one
    with $u=-d$ and another with the usual Standard Model
    hypercharges up to overall rescaling.
\end{enumerate}
\begin{solution}
    \begin{enumerate}[label=\roman*)]
    \item We start with the $SU(3)^3$ gauge anomaly, which only involves the quarks:
    \begin{equation}
        \sum_\text{left} A(\mathbf{3}) = \sum_\text{right} A(\mathbf{3}),
    \end{equation}
    which is indeed true. For $SU(2)^3$, we simply have $A(\mathbf{2}) = 0$, and hence is always anomaly-free. For $U(1)^3$, we have
    \begin{equation}
        \sum_\text{left}Y_i^3 = \sum_\text{right}Y_i^3 \implies 6q^3 + 2l^3 = 3u^3+ 3d^3+x^3.
    \end{equation}
    Now, we turn to the mixed gauge anomalies. For $SU(2)^2U(1)$, which only involves left-handed particles, we have
    \begin{equation}
        \sum_\text{left}Y_iI(\mathbf{2}) = \sum_\text{right}Y_iI(\mathbf{2}) \implies 3q + l = 0.
    \end{equation}
    For $SU(3)^2U(1)$, we find
    \begin{equation}
        \sum_\text{left}Y_iI(\mathbf{3}) = \sum_\text{right}Y_iI(\mathbf{3}) \implies 2q = u+d.
    \end{equation}
    Finally, $\text{grav}^2U(1)$ gives
    \begin{equation}
        \sum_\text{left} Y_i = \sum_\text{right} Y_i \implies 6q + 2l = 3u + 3d + x.
    \end{equation}
    \item The 4 equations boil down to
    \begin{equation}
        6u^3 +6d^3 + 21u^2d+21ud^2 = 0.
    \end{equation}
    Clearly, this has the solution $u/d = -1 \equiv y$. The polynomial factorises to
    \begin{equation}
        (y+1)(6y^2+15y+6) = 0,
    \end{equation}
    which has solutions, $y= -1, -2, -1/2$. Since $u_R$ and $d_R$ have the same representation, we can always relabel them. Therefore, we actually only have two possible classes of solutions:
    \begin{equation}
        u=-d, \quad q=0,\quad l=0,\quad x=0,
    \end{equation}
    \begin{equation}
        u=-2d,\quad q=\frac{d}{2},\quad l=\frac{3d}{2},\quad x =3d.
    \end{equation}
    \end{enumerate}
\end{solution}

\subsection{Sheet 3, Problem 4: Integer Hypercharges}

Repeat Problem 3 with $q,l,u,d,x\in\mathbb Z$.
\begin{enumerate}[label=\roman*)]
    \item Impose all gauge anomaly conditions, without imposing the
    mixed gravitational condition.
    \item Show that the solution is unique up to overall scaling and
    automatically cancels the mixed gravitational anomaly. You may
    quote a familiar pure-mathematics result without proving it.
\end{enumerate}
\textit{Hint:} First write $u-d=2y$ for integer $y$. Then use
the variables $v,w$ related by
\begin{equation}
    x=-\frac6{v+w},
    \qquad y=\frac{3(v-w)}{v+w}.
\end{equation}
\begin{solution}
    Since $u+d = 2q \in 2\mathbb{Z}$, we must have $u-d = 2y \in 2\mathbb{Z}$. Substituting this and $3q = -l$ into the $U(1)^3$ anomaly-free condition, we find
    \begin{equation}
        48q^3 +3(q+y)^3+3(q-y)^3 +x^3 = 0,
    \end{equation}
    which simplifies to
    \begin{equation}
        54q^3 +18qy^2 +x^3 = 0.
    \end{equation}
    Absorbing $q$ into $x,y$, so that $x,y \in \mathbb{Q}$ and
    \begin{equation}
        54 +18y^2 +x^3 = 0.
    \end{equation}
    Making the redefinition of variables and working through a bit of algebra, we find
    \begin{equation}
        v^3 +w^3 =1,
    \end{equation}
    which is the famous Fermat's last theorem with $N=3$. This, famously enough, has only trivial solutions, namely $v = 1, w=0$ or $v=0, w=1$. Again, these are equivalent solutions since we can always relabel $u$ and $d$. It is not hard to see that this is indeed the hypercharge of our world up to a scaling, and hence free of gravitational anomaly.
\end{solution}

\subsection{Sheet 3, Problem 5: Global Symmetries of the Standard Model}

Determine the global symmetries of the classical Standard Model
Lagrangian in each case:
\begin{enumerate}[label=\roman*)]
    \item there are no right-handed neutrinos;
    \item right-handed neutrinos exist and every allowed relevant and
    marginal interaction is included.
\end{enumerate}
For each case, reassess the result after accounting for anomalies.
\begin{solution}
    \begin{enumerate}[label=\roman*)]
    \item Naively, there seems to be a $U(1)^5$ global symmetry, which is
    \begin{align}
        U(1):\quad Q_L &\to e^{i\alpha}Q_L,\\
         d_R &\to e^{i\beta}d_R,\\
         u_R &\to e^{i\gamma}u_R,\\
         L_L &\to e^{i\delta}L_L,\\
         e_R &\to e^{i\lambda}e_R.
    \end{align}
        However, the Yukawa couplings impose additional restrictions on the charges, which are
    \begin{equation}
        \alpha = \beta, \quad \alpha = \gamma, \quad \delta = \lambda.
    \end{equation}
    This reduces the global symmetry to $U(1)^2$. In particular, we choose this to be $U(1)_B \times U(1)_L$.

    After accounting for the anomalies, the mixed anomaly with $SU(2)$ and $SU(3)$ vanishes for both lepton and baryon number symmetry. However, the mixed anomaly with hypercharge survives:
    \begin{equation}
        \mathcal{A}_{U(1)_BU(1)^2} = -\frac{1}{2} = -\mathcal{A}_{U(1)_LU(1)^2}.
    \end{equation}
    Therefore, the only surviving global symmetry in the quantum theory is $U(1)_B - U(1)_L$.
        \item The situation is almost identical after including right-handed neutrinos, since they have no hypercharge.
    \end{enumerate}

\end{solution}

\subsection{Sheet 3, Problem 6: Surviving Discrete Symmetries}

\textit{Part (a) is optional on the source sheet.}
\begin{enumerate}[label=\alph*)]
    \item Couple $SU(N)$ Yang--Mills theory to a single massless
    adjoint left-handed Weyl fermion $\lambda$. Explain quantum
    consistency. Although the theory has $\mathcal N=1$
    supersymmetry, that fact is not needed here. Classically
    $\lambda\mapsto e^{i\alpha}\lambda$ is a $U(1)$ symmetry.
    Follow the shift of the theta angle to show that only
    $\mathbb Z_{2N}$ survives quantum mechanically. If confinement
    produces $\langle\operatorname{Tr}\lambda\lambda\rangle
    \sim\Lambda_{\rm QCD}^3$, determine the action of
    $\mathbb Z_{2N}$ on this condensate and count the vacua.

    \item Now couple $SU(N)$, with $N>2$, to one Dirac fermion in
    the two-index symmetric representation. Determine its classical
    and quantum global symmetries. Assuming confinement and a
    condensate $\langle\bar\psi\psi\rangle
    \sim\Lambda_{\rm QCD}^3$, count the vacua. Repeat for the
    two-index antisymmetric representation. The representation
    data in Problem 2 may help.
\end{enumerate}

\begin{solution}
    \begin{enumerate}[label=\alph*)]
    \item Since $A(\mathrm{adj}) =0$, the theory is free of gauge anomaly, and hence consistent. For the $U(1)$ global symmetry, the ABJ anomaly is
    \begin{equation}
        \mathcal{A}_{U(1)SU(N)^2} = I(\mathrm{adj}) = 2N.
    \end{equation}
    Hence, the theta term transform as
    \begin{equation}
        \theta \to \theta + 2\alpha N.
    \end{equation}
    Therefore, only a discrete $\mathbb{Z}_{2N}$ survives. Under the discrete symmetry, the condensate transforms as
    \begin{equation}
        \langle\operatorname{Tr}\lambda\lambda\rangle \to e^{2in\pi/N}\langle\operatorname{Tr}\lambda\lambda\rangle.
    \end{equation}
    Therefore, the symmetry is further broken to $\mathbb{Z}_{2}$, which means that there are $N$ vacuum states for this theory.

    \item In this case the classical global symmetry is $U(1)_V \times U(1)_A$, while the quantum global symmetry reduces to $U(1)_V$ and a surviving discrete group of $U(1)_A$. In particular, the theta angle transform as
    \begin{equation}
        \theta \to \theta +2\beta(N\pm 2)
    \end{equation}
    ($+$ for symmetric representation and $-$ for antisymmetric representation). Therefore, the surviving discrete symmetry is $\mathbb{Z}_{N\pm 2}$. Under the discrete symmetry, the condensate transforms as
    \begin{equation}
        \langle\bar\psi\psi\rangle \to e^{2\pi ni/(N\pm 2)}\langle\bar\psi\psi\rangle.
    \end{equation}
    Therefore, there are $N\pm 2$ vacuum for the theory.
\end{enumerate}
\end{solution}

